%------------------------------------------------------------------------------
%	REQUIRED PACKAGES AND  CONFIGURATIONS
%------------------------------------------------------------------------------
%%%
%%%   DRACULA THEME
%%%
%\usepackage{/home/lorenzo/latex_dracula/draculatheme}
\documentclass{article}
%PACKAGES FOR TITLES
\usepackage{titlesec}
\usepackage{color}
\usepackage{tabularx}
\usepackage{multirow}
\usepackage[table]{xcolor}
\usepackage{titling}
\usepackage{nomencl}
\usepackage{biblatex} %Imports biblatex package
\addbibresource{bibliography.bib} %Import the bibliography file

% PACKAGES FOR LANGUAGE AND FONT
\usepackage[utf8]{inputenc}
\usepackage[english]{babel}
\usepackage[T1]{fontenc} % Font encoding

%TIKZ 
\usepackage{tikz}

% TIKZ PACKAGES
\usetikzlibrary{calc,patterns,angles,quotes}
\usepackage{pgfplots}
\pgfplotsset{compat = newest,width=6cm}
\usetikzlibrary{patterns.meta,decorations.pathmorphing}

% PACKAGES FOR IMAGES
\usepackage{graphicx}
\usepackage{float}
\usepackage{wrapfig}
\usepackage{import}
\usepackage{cancel}
\usepackage{xifthen}
\usepackage{pdfpages}
\usepackage{xcolor}
\usepackage{transparent}

\usepackage[acronym]{glossaries}

\makeglossaries

\newcommand{\incfig}[2]{%
    \ifthenelse{\isempty{#2}}{%
        \def\svgwidth{.75\columnwidth}
    }{%
        \def\svgwidth{#2\columnwidth}
    }%
    \import{./figures/}{#1.pdf_tex}
}

% PACKAGES FOR MATH
\usepackage{amsmath}
%\usepackage{siunitx}
% PACKAGES FOR REFERENCES & BIBLIOGRAPHY
\usepackage[colorlinks=true,linkcolor=black!40!black,anchorcolor=black,citecolor=black,filecolor=black,menucolor=black,runcolor=black,urlcolor=black]{hyperref} % Adds clickable links at references
\usepackage[english]{cleveref}
\pdfsuppresswarningpagegroup=1
%-------------------------------------------------------------------------
%	NEW COMMANDS DEFINED
%-------------------------------------------------------------------------
% EXAMPLES OF NEW COMMANDS -> here you see how to define new commands
\newcommand{\e}[1]{\times{} 10^{#1}}  % Powers of 10 notation
\newcommand{\mathbbm}[1]{\text{\usefont{U}{bbm}{m}{n}#1}} % From mathbbm.sty
\newcommand{\pdev}[2]{\frac{\partial#1}{\partial#2}}
\newcommand{\question}[1]{\colorbox{yellow}{#1}}
\newcommand{\integral}[4]{\int_{#1}^{#2} #3 \,\,\diff{} #4}
\newcommand{\dermat}[2]{\frac{\operatorname{D}#1}{\operatorname{D}#2}}
\newcommand{\abs}[1]{|#1|}
\newcommand{\person}[1]{\textsf{#1}}
\newcommand{\R}{\mathbbm{R}}
\newcommand{\costante}{\operatorname{cost}}
%\newcommand{\ln}{\operatorname{ln}}
\newcommand{\diff}{\operatorname{d}}
\newcommand{\vvline}{||}
\newcommand{\prodc}{\displaystyle\prod}
\newcommand{\parallelsum}{\mathbin{/\mkern-5mu/}}
\newcommand{\twomini}[2]{\begin{figure}[H]
                \centering
		\begin{minipage}{.5\linewidth}
			#1
		\end{minipage}
		\begin{minipage}{.5\linewidth}
			#2
		\end{minipage}
\end{figure}}
\newcommand{\twominisw}[4]{\begin{figure}[H]
                \centering
		\begin{minipage}{#3\linewidth}
			#1
		\end{minipage}
		\begin{minipage}{#4\linewidth}
			#2
		\end{minipage}
              \end{figure}}
            \newcommand{\threemini}[3]{\begin{figure}[H]
                \centering
                \begin{minipage}{.3\linewidth}
			#1
		\end{minipage}
		\begin{minipage}{.3\linewidth}
			#2
		\end{minipage}
		\begin{minipage}{.3\linewidth}
			#3
		\end{minipage}
\end{figure}}
%----------------------------------------------------------------------------
%	COLOURS 
%----------------------------------------------------------------------------
\definecolor{darkorange}{RGB}{229, 112, 0}
\definecolor{darkpurple}{RGB}{119, 25, 170}
\usepackage{xifthen}
\makeatother
\def\@lecture{}%
\newcommand{\lecture}[3]{
    \ifthenelse{\isempty{#3}}{%
        \def\@lecture{Lezione #1}%
    }{%
        \def\@lecture{Lezione #1: #3}%
    }%
    \subsection*{\@lecture}
    \marginpar{\small\textsf{\mbox{#2}}}
}

\usepackage{fancyhdr}
%\pagestyle{fancy}
\fancypagestyle{plain}{%
\fancyhf{} % cancella tutti i campi di intestazione e pi\‘e di pagina
\fancyfoot[C]{\bfseries \thepage} % tranne il centro
\renewcommand{\headrulewidth}{0pt}
\renewcommand{\footrulewidth}{0pt}}


% LE: left even
% RO: right odd
% CE, CO: center even, center odd
% My name for when I print my lecture notes to use for an open book exam.
% \fancyhead[LE,RO]{Lorenzo Pasqui}

% \fancyhead[RO,LE]{\@lecture} % Right odd,  Left even

% Right even, Left odd
%\fancyhead[R]{qualcosa}% TODO Aggiungi nomi 

% \fancyfoot[RO,LE]{}  % Right odd,  Left even
% \fancyfoot[RE,LO]{}          % Right even, Left odd
%\fancyfoot[C]{\thepage}     % Center

\makeatother

% redefinition of \maketitle with a logo ==============  
\makeatletter
\newcommand{\logo}[1]{\gdef\@logo{#1}}%
% \def\maketitle{
% \begin{center}
% \setlength\baselineskip{8ex}
% \setlength\parskip{4em plus 1fil minus 3em}
% \includegraphics[width=\textwidth,
% height=.6\textwidth,keepaspectratio]{\@logo}\par
% {\Huge\sffamily\bfseries \@title}\par
% {\Large\scshape \@author}\par
% \end{center}
% \newpage}
% \makeatother
%  end of redefinition  ===============================  \pdfsuppresswarningpagegroup=1
\logo{PoliLogo.jpg}
\author{Palusci Luca
	\\ 
	10730959 
\and
	 Pasqui Lorenzo 
	\\
        10703226
	\\
 \and
	Stellacci Andrea 
	\\ 
	10724258 
	\\
 \and
	Sharmoukh Carlo 
	\\ 
	Codice persona:
 }
 
	\date{11 Agosto 2024}
\newcommand{\myimage}{\includegraphics[width=0.2\textwidth]{relazione_latex/figures/PoliLogo.jpg}}
\newcommand{\mytitlepage}[1]{%
   \begin{center}
       \myimage
       \vspace{1cm} % Aggiungi spazio tra l'immagine e il titolo
       \LARGE #1
   \end{center}
}

%%% Local Variables:
%%% mode: latex
%%% TeX-master: f
%%% End:

\title{Appunti Tank}
\begin{document}
    \maketitle
    \tableofcontents
    \newpage
    \section{Volume del Tank}
    Il volume del tank deve essere maggiore di 
    quello del volume del propellente richiesto.

    Il tank è composto da 4 componenti:
    \begin{enumerate}
        \item $V_{pu}$: volume propellente (extra per emergenze).
        \item $V_{ull}$: volume \it{libero} per espansione del propellente o contrazione strutturale (.01/0.03).
        \item $V_{bo}$: per criogenici, permette l'ebollizione a causa del rifonrnimento e dello scolo.
        \item $V_{trap}$: volume del propellente che rimane intrappolato nelle feed lines (tipicamente volume del feed system).  
    \end{enumerate}
    Dunque: 
    \[V_{tot} = V_{pu} + V_{ull} + V_{bo} + V_{trap}\]
    \section{Forma del Tank}
    Di solito sono: 
    \begin{itemize}
        \item Sferici
        \item Cilindrici
    \end{itemize}
    I primi massimizzano il volume rispetto all'area, ma i cilindrici hanno forma 
    migliore per razzi e migliorano la rigidità strutturale.
    
    La pressione ha il maggior effetto sui limiti strutturali del tank. Per il design
    preliminare consideriamo solo i carichi di pressione:

    La design burst pressure del tank è:
    \[p_b = f_s \text{MEOP}\]
    con: 
    \begin{itemize}
        \item $p_b$: design bust pressure (in Pa).
        \item $f_s$: fattore sicurezza (2.0 per veicoli pressurizzati)
        \item $\text{MEPP}$: Maximum Expected Operating Pressure del tank (Pa).
    \end{itemize}
    \question{Quale materiale? serve densità ($\rho$), sforzo elastico ($F_{tu}$).}
    Indice di merito: \[\frac{F_{tu}}{\rho g_0}\]
    Altre cose da considerare nel materiale sono interazione 
    chimica e lavorazione (controllare in appendice B).
    \subsection{Sezione cilindrica}
    \begin{eqnarray*}
        t_s = \frac{p_b r_s}{2 F_{all}}\\
        m_s = A_s t_s \rho_{mat}
    \end{eqnarray*}
    con: 
    \begin{itemize}
        \item $r_s$ raggio della sfera
        \item $A_s$ superficie della sfera
        \item $V_s$ volume della sfera
        \item $t_s$ spessore
        \item $p_b$ DBP
        \item $F_{all}$ sforzo elastico permesso
        \item $m_s$ massa
        \item $\rho_{mat}$ densità del materiale
    \end{itemize}
    \subsection{Sezione sferica}
    \begin{eqnarray*}
        t_c = \frac{p_b r_c}{F_{all}}\\
        m_c = A_c t_c \rho_{mat}
    \end{eqnarray*}
    \subsubsection{Parti finali cilindriche}
    Non abbiamo ellissi perchè lo sforzo sul punto di contatto tra 
    ellissi e cilindro si formano sforzi alti. Abbiamo calotte sferiche.
    Non considerando i diversi loads dati da altri fattori che non siano la 
    pressione interna sottostimiamo di circa 2/2.5 volte.
    \section{Stima della massa con il metodo pV/W}
    Approccio puramente empirico. Considero $\phi_{tank}$, fattore della massa del tank:
    \begin{equation*}
        \phi_{tank} = \frac{p_b V_{tot}}{g_0 m_{tank}}
    \end{equation*}
    \emph{Per tank completamente metallici questo valore è 2500 metri!!}
    Posso quindi risolvere per la massa del tank.
    \subsection{Isolazione termica}
    Devo isolare liquidi criogenici o propellenti che potrebbero congelarsi in orbita. 
    Di solito questo è fatto tramite una lamina di metallo che copre una schiuma isolante (oppure materiale non metallico a nido d'ape).
    \subsection{Dispositivi di espulsione di cariburante}
    I serbatoi devono essere in grado di fornire e controllare il propellente in tutte le fasi della missione. 
    Parte di questa capacità deriva anche dalla rimozione di gas dal propellente che viene fornito al sistema propulsivo, e 
    espellere parte del propellente per eliminare i residui. Oltre a questo devono essere evitate forze trasmesse dal propellente alla navicella
    (slosh). Devo evitare che negli inlet finisca gas. 
    
    Posso avere: 
    \begin{itemize}
        \item sistemi di espulsione passivi
        \item sistemi di espulsione positivi
    \end{itemize}
    I primi meno costosi ed efficienti sfruttano la tensione superficiale del liquido, e sono dunque meno affidabili ma meno costosi e complicati, mentre i secondi usano barriere fisiche, valvole, pistoni, diaframmi etc.\ e sono quindi più efficaci a scapito di costo e complessità.
    \question{PAGINA 299 e precedenti}
    Posso inoltre usare gas per pressurizzare 
    \end{document}